\chapter{Frontend Design and Integration}
\section{Routing and Page Structure}

The frontend defines multiple entry routes to support desktop, mobile, embedded, and auxiliary usage scenarios. Each route maps to a dedicated UI layout optimised for its target environment:

\begin{itemize}
    \item \textbf{/} — Desktop chat interface
    \item \textbf{/mobile} — Mobile-optimised chat interface
    \item \textbf{/settings} — Configuration and settings page
    \item \textbf{/embed} — Full-page embedded interface
    \item \textbf{/ball} — Floating entry widget for embedding in third-party pages
    \item \textbf{/binarize.html} — Standalone static utility page
\end{itemize}

The root application component (\texttt{App.vue}) dynamically switches between desktop and mobile layouts according to browser viewport width. Embedded, floating, and settings views are excluded from this automatic switching to preserve their specific layout constraints.

\section{Backend Integration}

The backend service is started with:

\begin{verbatim}
uv run sustech-rag serve --host 127.0.0.1 --port 8001
\end{verbatim}

The frontend is launched separately with \texttt{npm run dev}. By default it routes all \texttt{/api} requests to the backend via a Vite proxy, so no absolute URL configuration is needed for local development. For cross-origin deployments the API base URL can be set to a full endpoint such as \texttt{http://127.0.0.1:8001/api}, and the backend CORS policy must be extended to permit the frontend's origin.

\section{State Management and Persistence}

Client-side state is persisted in browser \texttt{localStorage} under two keys:

\begin{itemize}
    \item \texttt{ragwebui:chats:v1} — conversation history across sessions
    \item \texttt{ragwebui:settings:v1} — user preferences and UI configuration
\end{itemize}

On initial load the frontend calls \texttt{POST /api/identity} to obtain a server-issued user identity. If the backend is unreachable, a locally generated fallback ID is used so the UI remains functional in offline or degraded conditions.

\section{SSE Streaming Protocol}

Model responses are delivered over a structured Server-Sent Events (SSE) stream. The protocol defines the following event types:

\begin{itemize}
    \item \texttt{start} — signals the beginning of a response stream
    \item \texttt{reference} — carries retrieved source chunks shown as citations
    \item \texttt{think.delta} / \texttt{think.end} — incremental chain-of-thought tokens and end marker
    \item \texttt{content.delta} — incremental response tokens
    \item \texttt{tool.call} / \texttt{tool.result} — reserved for agentic tool-use extensions
    \item \texttt{image} — reserved for multimodal output
    \item \texttt{error} — communicates recoverable or terminal errors to the client
    \item \texttt{done} — signals successful stream completion
\end{itemize}

The current backend emits \texttt{start}, \texttt{reference}, \texttt{think.delta}, \texttt{think.end}, \texttt{content.delta}, \texttt{error}, and \texttt{done}. A full protocol specification is maintained in \texttt{docs/API.md}.
